% Generated draft: requires mathematical and source review.
\documentclass[11pt]{article}
\usepackage[margin=25mm]{geometry}
\usepackage{fontspec}
\setmainfont{Latin Modern Roman}
\usepackage{amsmath,amssymb,mathtools,unicode-math}
\setmathfont{Latin Modern Math}
\usepackage{enumitem,needspace,xurl,hyperref}
\setlength{\parindent}{0pt}
\setlength{\parskip}{5pt}
\setlength{\emergencystretch}{3em}
\newcommand{\N}{\mathbb{N}}
\newcommand{\Z}{\mathbb{Z}}
\newcommand{\Q}{\mathbb{Q}}
\newcommand{\R}{\mathbb{R}}
\newcommand{\C}{\mathbb{C}}
\newcommand{\F}{\mathbb{F}}
\newcommand{\A}{\mathbb{A}}
\newcommand{\PP}{\mathbb{P}}
\newcommand{\E}{\mathbb{E}}
\newcommand{\eps}{\varepsilon}
\newcommand{\dd}{\,\mathrm{d}}
\begin{document}
\section*{UnsolvedMath ID: 30004211}
\textbf{Source problem code:} catalogue entry
\subsection{Definitions and mathematical statement}
Let \(A\) be a separable, unital, simple, exact \(C^*\)-algebra. Let \(T(A)\) denote its space of tracial states, equipped with the weak-* topology: these are positive linear functionals \(\tau:A\to\C\) satisfying \(\tau(1_A)=1\) and \(\tau(xy)=\tau(yx)\) for all \(x,y\in A\). Let \(\mathcal K\) be the compact operators on \(\ell^2(\N)\), and use the minimal tensor product \(A\otimes\mathcal K\). Its positive cone is denoted by \((A\otimes\mathcal K)^+\).

For positive elements \(a,b\in A\otimes\mathcal K\), write \(a\precsim b\) if there is a sequence \((v_n)\) in \(A\otimes\mathcal K\) such that \(v_n b v_n^*\to a\) in norm. Cuntz equivalence means comparison in both directions, and \([a]\) denotes the Cuntz equivalence class of \(a\). For \(\tau\in T(A)\), let \(\tau\otimes\operatorname{Tr}\) be the extended positive trace on \(A\otimes\mathcal K\), where \(\operatorname{Tr}\) is the usual unnormalized trace on \(\mathcal K\). Define
\[
 d_\tau([0])=0,\qquad
 d_\tau([a])=\lim_{n\to\infty}(\tau\otimes\operatorname{Tr})\!\left((a/\lVert a\rVert)^{1/n}\right)\in[0,\infty]
 \quad(a\ne0).
\]

The source uses \(\operatorname{LAff}(T(A))_{++}\) without specifying whether functions must be finite-valued. To make this unresolved convention explicit, for \(I=(0,\infty)\) or \(I=(0,\infty]\), let
\[
 \mathcal L_I(T(A))=\{f:T(A)\to I:\ f\text{ is lower semicontinuous and affine}\}.
\]
Here lower semicontinuity means that \(\{\tau:f(\tau)>r\}\) is open for every \(r\in\R\), and affinity means
\[
 f(t\tau+(1-t)\sigma)=t f(\tau)+(1-t)f(\sigma)
 \quad(\tau,\sigma\in T(A),\ 0<t<1),
\]
using extended-real arithmetic when \(I=(0,\infty]\). Thus strict positivity means positivity at every trace. The finite-valued and extended-valued interpretations are recorded separately; the supplied record does not determine which is intended.

\textbf{Source question.} For the intended choice of \(I\), is it true that for every separable, unital, simple, exact \(C^*\)-algebra \(A\) and every \(f\in\mathcal L_I(T(A))\), there exists \(a\in(A\otimes\mathcal K)^+\) such that
\[
 \forall\tau\in T(A),\qquad d_\tau([a])=f(\tau)?
\]
The source imposes no explicit infinite-dimensionality, stable-finiteness, or nonempty-trace-space hypothesis. These conditions have not been added. If \(T(A)=\varnothing\), the equality quantified over traces is vacuous. Whether any additional restriction on \(A\) was intended requires review.
\subsection{Short English statement}
The question asks whether every strictly positive lower semicontinuous affine function on the tracial state space of a separable, unital, simple, exact \(C^*\)-algebra occurs as the trace-rank function of a positive element of its stabilization. The source leaves the allowance of infinite function values and possible additional hypotheses on the algebra unresolved. Its open status is supplied dated curation, not independently verified.
\subsection{Sources}
\begin{itemize}
\item Supplied local UnsolvedMath registry record, permanent ID \(30004211\), title ``Realizing Lower Semicontinuous Trace Functions by Positive Elements,'' exact source problem code \texttt{catalogue entry}, in the Oberwolfach Reports Open Problems collection. No individual proposer is identified; the supplied proposed year is 2019.
\item Supplied report citation: \emph{C*-Algebras (2020)}, Oberwolfach Reports, DOI \texttt{10.4171/owr/2019/37}, pp.\ 2257--2332. \url{https://doi.org/10.4171/owr/2019/37}. The literature-source list also labels this reference ``Oberwolfach Report 37, Cuntz semigroup.'' The citation's 2020 designation and the record's 2019 proposed year are retained as supplied.
\item Supplied literature reference: Dadarlat and Toms, ``ranks of operators,'' DOI \texttt{10.1016/j.jfa.2010.03.022}. \url{https://doi.org/10.1016/j.jfa.2010.03.022}. This is the abbreviated attribution supplied by the record.
\item Literature assessment dated exactly \texttt{2026-08-22}: the supplied curation labels the problem open and reports rank realization in regular classes, injectivity under strict comparison, and an unresolved general surjectivity problem. This is supplied dated curation derived from the OWR record, not an independently verified assessment. The scope of the injectivity claim and the meaning of ``regular classes'' are not specified.
\item The registry labels the statement \texttt{corrected\_verified} and says its mathematical content was preserved while notation, typography, duplicate clauses, or omitted definitions were repaired. These are supplied provenance claims; neither the cited sources nor mathematical correctness have been independently verified here.
\end{itemize}
\end{document}
